\documentclass{article}
\usepackage[T1]{fontenc}
\usepackage{lmodern}
\usepackage{graphicx}
\usepackage{amsmath,amssymb}
\usepackage[a4paper,margin=2cm]{geometry}
\usepackage[absolute,overlay]{textpos}
\setlength{\TPHorizModule}{1mm}
\setlength{\TPVertModule}{1mm}
\usepackage[indonesian]{babel}
\usepackage{hyperref}
\setlength{\emergencystretch}{2em}
% unit: Halaman 6

\begin{document}

% === Halaman 6 (llm) ===

\begin{align*}
a^{\phi(n)} &\equiv 1 \pmod{n} \\
&\downarrow \text{ (pangkatkan kedua ruas dengan k)} \\
a^{k\phi(n)} &\equiv 1^k \pmod{n} \\
&\downarrow \\
a^{k\phi(n)} &\equiv 1 \pmod{n} \\
&\downarrow \text{ (ganti $a$ dengan $m$)} \\
m^{k\phi(n)} &\equiv 1 \pmod{n} \\
&\downarrow \text{ (kalikan kedua ruas dengan $m$)} \\
m^{k\phi(n)+1} &\equiv m \pmod{n}
\end{align*}

\end{document}
